% Isolated CP35 candidate; no canonical or shared build input is changed.
\subsection{Reduced twisted cross determinants on both product axes}

Let \(k\) be algebraically closed of characteristic two,
\(E:y^2+y=x^3\), \(M_E=\mathcal O_E(3O)\), and
\(M=M_E\boxtimes M_E\) on \(E^2\).
Let \(P=P_1\boxtimes P_2\), with both \(P_i\ne0\).
Write \(P_1=\mathcal O_E((a,b)-O)\) and
\(P_2=\mathcal O_E((c,d)-O)\), so
\(b^2+b=a^3\) and \(d^2+d=c^3\).
All these points are finite, since the omitted origin is the sole
point at infinity. The matrices below use the actual relative
Cartier bases of Lemma~\ref{lem:DE-SS-cross-elliptic}.

\begin{lemma}[Two fixed lines certify every two-axis twist]
\label{lem:ss-product-two-axis-universal-line}
\status{Candidate; universal determinant identity and direct proof supplied.}
Put
\[
 A_0(a,b)=\begin{pmatrix}1&a&0\\a&0&b\\0&b+1&a\end{pmatrix},
 \quad A_1=\begin{pmatrix}0&0&1\\0&1&0\\1&0&0\end{pmatrix},
 \quad A_2=\begin{pmatrix}0&0&0\\0&0&1\\0&1&0\end{pmatrix}.
 \tag{TAS1}
\]
For a \(3\times3\) functional coefficient matrix \(X\), set
\[
 C_{a,b,c,d}(X)=\sum_{i,j=0}^2X_{ij}A_i(a,b)\otimes A_j(c,d).
\]
On the projective line
\[
 X(T,U)=\begin{pmatrix}
 0&0&T\\U&U+T&U+T\\U+T&U&U+T
 \end{pmatrix},
 \tag{TAS2}
\]
its determinant is the following homogeneous binary form of degree nine:
\[
 \begin{split}
 F(T,U)={}&u^2T^9+v^2T^8U
              +w^2(T^7U^2+T^6U^3)+TU^8+U^9,\\
 u={}&b+d+ac+a+1,\qquad v=c(a+1),\qquad w=ac+a+c.
 \end{split}
 \tag{TAS3}
\]
If \(\Delta=b+d+a+c+1\ne0\), this binary form is squarefree.
If \(\Delta=0\), the determinant of the same original cross
pencil is squarefree on the other fixed line obtained from (TAS2)
by replacing its third row by the sum of its first and third rows.
\end{lemma}
\begin{proof}
Expansion of the actual nine-by-nine determinant on \(U=1\),
followed by the two monic curve relations, gives
\[
 \begin{split}
 &(a^3+a^2c^2+a^2+b+c^3+d+1)T^9
       +(a^2c^2+c^2)T^8\\
 &\hspace{15mm}+(a^2c^2+a^2+c^2)(T^7+T^6)+T+1.
 \end{split}
 \tag{TAS4}
\]
The first coefficient is \(u^2\), because
\(b^2=b+a^3\) and \(d^2=d+c^3\); the next coefficients
are \(v^2,w^2\). Homogenization proves (TAS3).
The accompanying standalone binary checker independently reconstructs
the full determinant by the subset permutation recurrence, reducing
only the two displayed monic curve equations. Thus (TAS4) is a
universal polynomial identity, not a specialization or sample.

The two homogeneous partial derivatives are
\[
 F_T=(uT^4+wT^3U+U^4)^2,\qquad
 F_U=(vT^4+wT^3U+U^4)^2.
 \tag{TAS5}
\]
Their square roots differ by
\((u+v)T^4=\Delta T^4\). If \(\Delta\ne0\), a common
projective zero would have \(T=0\) and then \(U=0\), which
is impossible. Thus the binary form is squarefree, including at
infinity. This argument allows \(u=0\); it does not incorrectly
discard the possible simple root at infinity when the affine
degree drops to eight.

If \(\Delta=0\), replace \(b\) by \(b+1\). This is the
inverse point of the first elliptic factor and changes \(\Delta\)
to one. Moreover
\[
                    A_0(a,b+1)=A_0(a,b)+A_2.
 \tag{TAS6}
\]
Consequently \(C_{a,b+1,c,d}(X)\) equals
\(C_{a,b,c,d}(X')\), where \(X'\) has third row equal
to the sum of the first and third rows of \(X\).
This is an invertible linear change on the full functional space,
and carries the line (TAS2) to the second stated line. The preceding
argument applied to \(b+1\) therefore proves a squarefree
restriction for the original twist. No change of its order or
an assumption about distribution of torsion points is used.
\end{proof}

\begin{theorem}[Uniform two-axis reducedness and generic corank one]
\label{thm:ss-product-two-axis-cross-squarefree}
\status{Candidate; complete deduction supplied.}
For every pair of nontrivial degree-zero line bundles \(P_1,P_2\)
on \(E\), the actual nine-variable relative Cartier cross
determinant for \(P_1\boxtimes P_2\) is geometrically squarefree.
Every irreducible determinant component has generic cross corank one.
In particular these assertions hold for arbitrary nonzero odd-order
torsion on both factors, with no bound on the odd order.
\end{theorem}
\begin{proof}
The preceding lemma provides one of two projective lines on which
the homogeneous degree-nine determinant restricts to a nonzero
squarefree binary form. If a homogeneous irreducible factor of the
full determinant occurred with multiplicity at least two, its
restriction would be a nonzero homogeneous binary form of positive
degree: it cannot vanish identically because the determinant
restriction is nonzero. Over \(k\) it has a projective zero,
which would be a repeated zero of the determinant restriction.
This is a contradiction. Thus the full determinant is squarefree.

At the generic point of a reduced determinant component some first
derivative is nonzero. Corank at least two of a matrix of linear
forms would make every eight-dimensional cofactor zero there,
and hence every derivative of its determinant zero. The generic
rank is therefore eight. The proof treats every nonzero \(P_i\),
so its torsion specialization requires no further existence or
openness argument.
\end{proof}

\begin{corollary}[What persists in an actual deformation]
\label{cor:ss-product-two-axis-reduced-deformation-open}
\status{Candidate; local openness only.}
Suppose a family of polarized abelian surfaces with a degree-zero
twist specializes to the above product data. On a neighborhood
where the three nine-dimensional spaces in the relative cross
pairing are locally free, geometric reducedness of the cross
determinant holds on a nonempty open containing that product point.
The twist may be a specified odd-torsion section of the family.
\end{corollary}
\begin{proof}
Trivialize the input and target bundles locally, with bases
specializing to those above. The two fixed functional directions
of the appropriate line lift to the coefficient bundle. They
remain independent after shrinking. The common-zero scheme of
the two homogeneous derivatives of the restricted binary form
is closed in the relative projective line, hence has closed image
in the base by properness. It has empty fibre at the product
point by (TAS5). Its complement is the required open.
The preceding factor-restriction argument applies in every
geometric fibre in this open. This proves only a statement within
the given family; it does not assert that every fixed surface
or every chosen odd-torsion section meets such a product point.
\end{proof}

\begin{remark}[Cycle and global-image scope]
\label{rem:ss-product-two-axis-squarefree-scope}
\status{Candidate scope; no new full-monodromy conclusion.}
The one-axis determinant-cube example and the present two-axis
theorem distinguish actual product twists. A reduced cross pencil
is not yet a reduced determinant of a composed Frobenius-cycle
operator. Neither stable corank one, full inertia over a character
field, irreducibility of a cycle local system, nor full monodromy
is asserted here. The deformation corollary makes no claim that
all fixed nonproduct surfaces are reached with their prescribed
torsion data.

The discovery computations are retained separately, including the
explicit order-three example whose line polynomial is \(T^9+T+1\).
The standalone reconstruction
\path{code/monodromy/verify_ss_product_two_axis_squarefree.py}
uses only binary polynomial arithmetic, the two monic elliptic
relations and three independent subset determinant reconstructions.
It checks (TAS3), both affine forms of (TAS5), (TAS6) on the
changed line, and the exact change of \(\Delta\) by one.
Its author run passed at 04:38:32 UTC on 20 November 2026 in
\(0.03699702300946228\) seconds. No finite-parameter sampling
is used for the universal theorem.
\end{remark}
